% preamble.tex - 中文线性代数讲义
% XeLaTeX 配置 (ctex 中文支持)
% 编译方式: xelatex main.tex (运行 2 次以生成目录和引用)

\documentclass[12pt,a4paper,openany,oneside]{book}

% === 中文支持 (必须最先加载, 需要 XeLaTeX) ===
\usepackage[UTF8, scheme=plain]{ctex}

% === 页面设置 ===
\usepackage{geometry}
\geometry{
    top=2.5cm,
    bottom=2.5cm,
    left=2.5cm,
    right=2.5cm,
    headheight=14pt
}

% === 数学宏包 ===
\usepackage{amsmath}
\usepackage{amssymb}
\usepackage{amsthm}
\usepackage{mathtools}
\usepackage{bm}

% === 图形与绘图 ===
\usepackage{graphicx}
\usepackage{tikz}
\usepackage{pgfplots}
\pgfplotsset{compat=1.18}
\usetikzlibrary{arrows.meta, calc, positioning, decorations.pathreplacing, patterns}

% === 颜色与盒子 ===
\usepackage{xcolor}
\usepackage{tcolorbox}
\tcbuselibrary{skins,breakable,theorems}

% === 表格与列表 ===
\usepackage{booktabs}
\usepackage{array}
\usepackage{enumitem}

% === 超链接 (建议放在较后位置) ===
\usepackage[colorlinks=true,linkcolor=blue!70!black,urlcolor=blue!70!black,citecolor=green!50!black]{hyperref}

% === 定理环境 (中文) ===
\theoremstyle{definition}
\newtheorem{definition}{定义}[chapter]
\newtheorem{example}{例}[chapter]
\newtheorem{exercise}{练习}[chapter]
\newtheorem{problem}{习题}[chapter]

\theoremstyle{plain}
\newtheorem{theorem}{定理}[chapter]
\newtheorem{lemma}{引理}[chapter]
\newtheorem{proposition}{命题}[chapter]
\newtheorem{corollary}{推论}[chapter]

\theoremstyle{remark}
\newtheorem{remark}{注}[chapter]
\newtheorem{note}{备注}[chapter]

% === 自定义 tcolorbox 环境 (中文标题) ===

% 几何直观盒子
\newtcolorbox{intuition}{
    enhanced,
    colback=blue!5,
    colframe=blue!75!black,
    fonttitle=\bfseries,
    title={几何直观},
    breakable,
    left=5mm,
    right=5mm,
    top=3mm,
    bottom=3mm,
    arc=2mm,
    boxrule=1pt
}

% 实际应用盒子
\newtcolorbox{practical}{
    enhanced,
    colback=green!5,
    colframe=green!60!black,
    fonttitle=\bfseries,
    title={实际应用},
    breakable,
    left=5mm,
    right=5mm,
    top=3mm,
    bottom=3mm,
    arc=2mm,
    boxrule=1pt
}

% 注意盒子
\newtcolorbox{warning}{
    enhanced,
    colback=red!5,
    colframe=red!75!black,
    fonttitle=\bfseries,
    title={注意},
    breakable,
    left=5mm,
    right=5mm,
    top=3mm,
    bottom=3mm,
    arc=2mm,
    boxrule=1pt
}

% 重点盒子
\newtcolorbox{important}{
    enhanced,
    colback=orange!5,
    colframe=orange!75!black,
    fonttitle=\bfseries,
    title={重点},
    breakable,
    left=5mm,
    right=5mm,
    top=3mm,
    bottom=3mm,
    arc=2mm,
    boxrule=1pt
}

% 本章小结盒子
\newtcolorbox{summary}{
    enhanced,
    colback=purple!5,
    colframe=purple!75!black,
    fonttitle=\bfseries,
    title={本章小结},
    breakable,
    left=5mm,
    right=5mm,
    top=3mm,
    bottom=3mm,
    arc=2mm,
    boxrule=1pt
}

% === 页眉页脚 ===
\usepackage{fancyhdr}
\pagestyle{fancy}
\fancyhf{}
\fancyhead[L]{\leftmark}
\fancyhead[R]{\rightmark}
\fancyfoot[C]{\thepage}
\renewcommand{\headrulewidth}{0.4pt}
\renewcommand{\footrulewidth}{0pt}

% === 行距 ===
\usepackage{setspace}
\onehalfspacing

% === 摘要环境 (book 文档类默认没有 abstract 环境, 自定义为缩进小字块) ===
\newenvironment{abstract}{%
  \small
  \begin{quotation}\noindent\ignorespaces
}{%
  \end{quotation}
}

% === 自定义数学命令 ===
% commands.tex - 线性代数自定义数学命令
% 中文线性代数讲义 (与英文版一致, 另补英文版漏定义的 \vp 等)

% === Vector Notation ===
% Bold vectors (column vectors)
\newcommand{\vect}[1]{\mathbf{#1}}
\newcommand{\vv}{\vect{v}}
\newcommand{\vw}{\vect{w}}
\newcommand{\vu}{\vect{u}}
\newcommand{\vx}{\vect{x}}
\newcommand{\vy}{\vect{y}}
\newcommand{\vz}{\vect{z}}
\newcommand{\va}{\vect{a}}
\newcommand{\vb}{\vect{b}}
\newcommand{\vc}{\vect{c}}
\newcommand{\vd}{\vect{d}}
\newcommand{\ve}{\vect{e}}
\newcommand{\vf}{\vect{f}}
\newcommand{\vg}{\vect{g}}
\newcommand{\vh}{\vect{h}}
\newcommand{\vl}{\vect{l}}
\newcommand{\vm}{\vect{m}}
\newcommand{\vn}{\vect{n}}
\newcommand{\vo}{\vect{o}}
\newcommand{\vp}{\vect{p}}
\newcommand{\vq}{\vect{q}}
\newcommand{\vr}{\vect{r}}
\newcommand{\vs}{\vect{s}}
\newcommand{\vt}{\vect{t}}

% Basis vectors (with hat notation)
\newcommand{\vi}{\hat{\imath}}
\newcommand{\vj}{\hat{\jmath}}
\newcommand{\vk}{\hat{k}}

% Zero vector
\newcommand{\vzero}{\vect{0}}

% === Matrix Notation ===
\newcommand{\mat}[1]{\mathbf{#1}}
\newcommand{\mA}{\mat{A}}
\newcommand{\mB}{\mat{B}}
\newcommand{\mC}{\mat{C}}
\newcommand{\mD}{\mat{D}}
\newcommand{\mE}{\mat{E}}
\newcommand{\mF}{\mat{F}}
\newcommand{\mG}{\mat{G}}
\newcommand{\mH}{\mat{H}}
\newcommand{\mI}{\mat{I}}
\newcommand{\mJ}{\mat{J}}
\newcommand{\mK}{\mat{K}}
\newcommand{\mL}{\mat{L}}
\newcommand{\mM}{\mat{M}}
\newcommand{\mN}{\mat{N}}
\newcommand{\mO}{\mat{O}}
\newcommand{\mP}{\mat{P}}
\newcommand{\mQ}{\mat{Q}}
\newcommand{\mR}{\mat{R}}
\newcommand{\mS}{\mat{S}}
\newcommand{\mT}{\mat{T}}
\newcommand{\mU}{\mat{U}}
\newcommand{\mV}{\mat{V}}
\newcommand{\mW}{\mat{W}}
\newcommand{\mX}{\mat{X}}
\newcommand{\mY}{\mat{Y}}
\newcommand{\mZ}{\mat{Z}}

% Matrix with brackets
\newcommand{\bmat}[1]{\begin{bmatrix}#1\end{bmatrix}}
\newcommand{\pmat}[1]{\begin{pmatrix}#1\end{pmatrix}}
\newcommand{\vmat}[1]{\begin{vmatrix}#1\end{vmatrix}}

% === Common Operations ===
% Transpose
\newcommand{\trans}{^{\mathsf{T}}}

% Inverse
\newcommand{\inv}{^{-1}}

% Norm and absolute value
\newcommand{\norm}[1]{\left\| #1 \right\|}
\newcommand{\abs}[1]{\left| #1 \right|}

% === Important Spaces ===
\newcommand{\R}{\mathbb{R}}
\newcommand{\C}{\mathbb{C}}
\newcommand{\N}{\mathbb{N}}
\newcommand{\Z}{\mathbb{Z}}
\newcommand{\Q}{\mathbb{Q}}

% Rn space
\newcommand{\Rtwo}{\mathbb{R}^2}
\newcommand{\Rthree}{\mathbb{R}^3}

% === Linear Algebra Operators ===
\DeclareMathOperator{\spn}{span}
\DeclareMathOperator{\Span}{Span}
\DeclareMathOperator{\Null}{Null}
\DeclareMathOperator{\Col}{Col}
\DeclareMathOperator{\Row}{Row}
\DeclareMathOperator{\rank}{rank}
\DeclareMathOperator{\tr}{tr}
\DeclareMathOperator{\proj}{proj}

% === Common Matrices ===
% 2x2 matrix shorthand
\newcommand{\twomat}[4]{\begin{bmatrix} #1 & #2 \\ #3 & #4 \end{bmatrix}}

% 3x3 matrix shorthand
\newcommand{\threemat}[9]{\begin{bmatrix} #1 & #2 & #3 \\ #4 & #5 & #6 \\ #7 & #8 & #9 \end{bmatrix}}

% 2D column vector
\newcommand{\twovec}[2]{\begin{bmatrix} #1 \\ #2 \end{bmatrix}}

% 3D column vector
\newcommand{\threevec}[3]{\begin{bmatrix} #1 \\ #2 \\ #3 \end{bmatrix}}

% === Change of Basis ===
\newcommand{\coords}[2]{[#1]_{#2}}

% === System of Equations ===
\newcommand{\system}[1]{\left\{\begin{aligned} #1 \end{aligned}\right.}

% === Emphasizing ===
\newcommand{\vocab}[1]{\textbf{\textcolor{blue!70!black}{#1}}}


% === TikZ 线性代数绘图样式 ===
\tikzset{
    vector/.style={-{Stealth[length=3mm]}, thick},
    axis/.style={-{Stealth[length=2mm]}, gray},
    grid/.style={very thin, gray!30},
    transformed/.style={blue!70},
    original/.style={black},
    basis/.style={red!70!black, very thick},
    point/.style={circle, fill, inner sep=1.5pt}
}
